% titlepage.tex -- титульный лист.
% Полностью повторяет форму из zapiska/Tit_list_PP.pdf (бланк титульного листа
% отчёта по практике, кафедра ИСиТ СЛИ): шапка по центру, наименование работы
% полужирным, ниже -- сведения о работе и подписи, внизу листа -- "Сыктывкар".
% Рамка (как была в прежнем титульном листе пояснительной записки) бланком не
% предусмотрена, поэтому не рисуется.
%
% Подстановочные поля бланка (в PDF выделены заливкой) заполнены данными
% работы: тема и фамилия исполнителя.

% Линия для подписи и подпись под ней -- одной ширины, чтобы пояснение
% "(подпись, дата)" стояло ровно по центру линии.
\newcommand{\signline}{\rule{4.2cm}{0.5pt}}
\newcommand{\signcap}{\makebox[4.2cm][c]{\footnotesize (подпись, дата)}}

% ---------- Поля титульного листа ----------
% Титульный лист набран по бланку и правому полю 15 мм, принятому для
% основного текста, не подчиняется: графы подписей заданы в сантиметрах по
% бланку и рассчитаны на полосу 170 мм. Поэтому на время титульного листа
% правое поле остаётся прежним (10 мм), а сразу после него геометрия
% возвращается к основной.
\newgeometry{left=30mm,right=10mm,top=20mm,bottom=20mm}
\begin{titlepage}
\begin{spacing}{1.0}
    \centering
    % Шапка прижата к верхнему полю: отрицательный отступ съедает \topskip
    % первой строки, поэтому текст начинается почти от границы поля 20 мм.
    % Освободившееся место уходит в \vfill перед "Сыктывкар 2026",
    % так что нижняя строка остаётся на месте.
    \vspace*{-\baselineskip}

    МИНИСТЕРСТВО НАУКИ И ВЫСШЕГО ОБРАЗОВАНИЯ\\
    РОССИЙСКОЙ ФЕДЕРАЦИИ

    \vspace{\baselineskip}
    СЫКТЫВКАРСКИЙ ЛЕСНОЙ ИНСТИТУТ (ФИЛИАЛ) ФЕДЕРАЛЬНОГО\\
    ГОСУДАРСТВЕННОГО БЮДЖЕТНОГО ОБРАЗОВАТЕЛЬНОГО\\
    УЧРЕЖДЕНИЯ ВЫСШЕГО ОБРАЗОВАНИЯ\\
    «САНКТ-ПЕТЕРБУРГСКИЙ ГОСУДАРСТВЕННЫЙ ЛЕСОТЕХНИЧЕСКИЙ\\
    УНИВЕРСИТЕТ ИМЕНИ С.\,М.\,КИРОВА» (СЛИ)

    \vspace{\baselineskip}
    ТРАНСПОРТНО-ТЕХНОЛОГИЧЕСКИЙ ФАКУЛЬТЕТ

    \vspace{\baselineskip}
    КАФЕДРА ИНФОРМАЦИОННЫХ СИСТЕМ И ТЕХНОЛОГИЙ

    \vspace{\baselineskip}
    % 13 pt: при 14 pt первая строка бланка не умещается в ширину полосы
    % и разрывается не там, где в бланке.
    {\fontsize{14}{16}\selectfont\bfseries
    ОТЧЁТ ПО ТЕХНОЛОГИЧЕСКОЙ (ПРОЕКТНО-ТЕХНОЛОГИЧЕСКОЙ)\\
    ПРАКТИКЕ. WEB-ТЕХНОЛОГИИ\par}

    \vspace{\baselineskip}
    \raggedright
    на тему: Разработка CMS для парикмахерских и барбершопов «Сайтама»\\
    Место прохождения практики: СЛИ\\
    Период прохождения: 22.06.26 -- 18.07.26

    \vspace{\baselineskip}
    \begin{tabular}{@{}p{7.6cm}@{\hspace{0.3cm}}p{4.2cm}@{\hspace{0.3cm}}p{4.4cm}@{}}
        Выполнил: студент 3 курса     &            &                   \\
        очной формы обучения          &            &                   \\
        направления подготовки        & \signline  & Решетняк В.\,А.   \\
        бакалавриата «Информационные  & \signcap   &                   \\
        системы и технологии»         &            &                   \\
    \end{tabular}

    \vspace{\baselineskip}
    \begin{tabular}{@{}p{7.6cm}@{\hspace{0.3cm}}p{4.2cm}@{\hspace{0.3cm}}p{4.4cm}@{}}
        Руководитель:                 & \signline  & Оганезова Н.\,А.  \\
        к.\,э.\,н., доцент            & \signcap   &                   \\
    \end{tabular}

    \vspace{\baselineskip}
    \begin{tabular}{@{}p{7.6cm}@{\hspace{0.3cm}}p{4.2cm}@{\hspace{0.3cm}}p{4.4cm}@{}}
        Заведующий кафедрой           &            &                   \\
        информационных систем и       & \signline  & Плешев Д.\,А.     \\
        технологий, к.\,ф.-м.\,н., доцент & \signcap &                 \\
    \end{tabular}

    \vspace{1.5\baselineskip}
    Отчёт защищён на оценку\,\rule{6cm}{0.5pt}

    \vspace{0.5\baselineskip}
    «\rule{1cm}{0.5pt}» \rule{4.5cm}{0.5pt} 2026 г.

    \vfill
    \centering
    Сыктывкар 2026
\end{spacing}
\end{titlepage}
\restoregeometry
